\documentclass[11pt,a4paper]{article}

\usepackage[margin=1in]{geometry}
\usepackage{amsmath, amssymb}
\usepackage{graphicx}
\usepackage{hyperref}
\usepackage{listings}
\usepackage{siunitx}
\usepackage{booktabs}
\usepackage{longtable}
\usepackage{xcolor}
\usepackage{framed}

\lstset{
  basicstyle=\ttfamily\small,
  breaklines=true,
  frame=single,
  showstringspaces=false,
}

% Inline code. Usable in math mode, where \texttt alone would be ignored.
\newcommand{\fn}[1]{\ifmmode\text{\ttfamily #1}\else\texttt{#1}\fi}

% Callout for behaviour that departs from documented or expected intent.
\newenvironment{warningbox}%
  {\par\begin{framed}\noindent\textbf{Discrepancy.}\ \ignorespaces}%
  {\end{framed}\par}

\title{Doppler Ultrasound Segmentation and Metric Extraction Pipeline}
\author{Dale Kernot}
\date{\today}

\begin{document}
\maketitle
\tableofcontents

\newpage
\section{Introduction}
  \subsection{Motivation and clinical context}
    Doppler ultrasound scans contain two linked sources of information. The
    first is the image of the Doppler waveform, which can be segmented and
    digitised into a velocity-time trace. The second is the machine-generated
    text panel, which commonly reports clinical measurements such as peak
    systolic velocity, end diastolic velocity, heart rate, pulsatility index,
    and resistive index.\\
    \\
    The aim of this pipeline is to extract both sources of information from a
    scan, compare them where possible, and return a structured table of
    measurements together with a digitised waveform trace. The code is designed
    around scans with a clearly visible Doppler waveform and machine annotation,
    but the implementation contains fallback behaviour for incomplete axes,
    missing labels, and scans where only part of the expected visual structure is
    present.\\
    \\
    Two fundamentally different classes of input are supported, and they are
    calibrated in fundamentally different ways. Screenshot-style raster images
    (JPEG, PNG and similar) carry no calibration metadata, so the pipeline must
    recover the velocity scale by reading the axis tick marks and numeric labels
    off the image itself. DICOM files carry explicit spatial calibration in their
    metadata, so the pixel-to-physical-unit mapping is read directly rather than
    inferred. Keeping these two paths clearly separated is essential to
    understanding the codebase, and this document treats them as distinct
    branches throughout.

  \subsection{Scope of this document}
    This document describes the processing pipeline implemented in the
    \fn{usseg} package. It is organised by processing stage. For each stage it
    states what the stage is trying to achieve, which functions implement it,
    what data structures flow in and out, and how the raster and DICOM paths
    differ.\\
    \\
    The single-image entry point, \fn{data\_from\_image}, is used as the
    reference implementation of the core algorithm because it contains the
    algorithmic path with the least surrounding bookkeeping. The batch pipeline
    performs the same algorithmic stages and adds file discovery, output
    writing, annotation rendering, serialisation and report generation.

  \subsection{Document conventions}
    Function and module names are written in a monospaced font, for example
    \fn{segment\_refinement}. Unless stated otherwise, functions named without a
    module prefix live in \fn{general\_functions.py}, which holds the majority of
    the algorithmic code.\\
    \\
    Image coordinates follow the convention used in the code: arrays are indexed
    as \texttt{[row, column]}, equivalently \texttt{[y, x]}, with the row index
    increasing downwards from the top of the image. Consequently a numerically
    smaller row index is visually higher in the image, which matters when reading
    the envelope-orientation logic in Section~\ref{sec:segmentation}.

  \subsection{Codebase overview}
    The modules involved in the processing pipeline are:

    \begin{description}
      \item[\fn{single\_image\_processing.py}] Provides \fn{data\_from\_image},
        the single-scan API used by examples and tests.
      \item[\fn{segment\_files.py}] Provides \fn{segment}, the batch workflow
        over a list of scan files, including annotated and digitised output
        rendering and serialisation.
      \item[\fn{general\_functions.py}] Contains the majority of the
        implementation: OCR and text repair, coarse and refined segmentation, the
        three waveform tracing methods, axis detection, digitisation, beat
        detection, metric calculation and annotation.
      \item[\fn{hemodynamic\_indices.py}] Pure functions for the haemodynamic
        indices (pulsatility index, resistive index, time-averaged maximum
        velocity).
      \item[\fn{digitized\_comparison.py}] Compares digitised measurements
        against OCR-derived measurements and selects the best-matching
        segmentation model.
      \item[\fn{organise\_files.py}] Discovers candidate ultrasound files within
        a configured directory tree.
      \item[\fn{visualisation\_html.py}] Builds the HTML report used to inspect
        batch results.
      \item[\fn{setup\_environment.py}] Configures the Tesseract OCR
        environment.
      \item[\fn{main.py}] Batch entry point that reads \fn{config.toml},
        discovers files, segments them and generates the HTML report.
    \end{description}

\newpage
\section{Architecture and Entry Points}
  \subsection{Single-image entry point}
    \fn{data\_from\_image} in \fn{single\_image\_processing.py} processes exactly
    one scan and returns its results in memory without writing any files. It is
    the recommended way to inspect the behaviour of the algorithm on a single
    input.

    \begin{lstlisting}[language=Python]
from usseg import data_from_image

df, (xdata, ydata) = data_from_image(image_path="path/to/scan.png")
    \end{lstlisting}

    The function returns two objects:

    \begin{description}
      \item[\fn{df}] A pandas dataframe. For raster images this contains the
        OCR-extracted text metrics, the digitised equivalents of those metrics
        for each segmentation model, comparison fields, and a
        \texttt{Returned model} entry recording which model was selected. For
        DICOM input it contains metrics derived purely from the digitised
        waveform, optionally preceded by a vessel/side label row.
      \item[\fn{[xdata, ydata]}] The digitised waveform trace. For raster images
        this is the trace of the selected model after heart-rate time scaling.
        For DICOM it is the primary digitised series in physical units.
    \end{description}

    A legacy calling convention accepting pre-loaded Pillow and OpenCV image
    objects is still supported, but it cannot express the DICOM path and is
    scheduled for removal. Passing \fn{image\_path} together with either legacy
    argument raises \fn{ValueError}.

  \subsection{Batch entry point}
    The batch pipeline is driven by \fn{main.py}, which reads \fn{root\_dir}
    from \fn{config.toml} and calls, in order, \fn{setup\_tesseract},
    \fn{get\_likely\_us}, \fn{segment} and \fn{generate\_html\_from\_pkl}. Each
    of these three processing calls is wrapped in the \fn{prof} helper, which
    profiles the call with \fn{cProfile} and writes cumulative-time statistics to
    a per-function log file. Running the module directly also configures logging
    at \texttt{INFO} level into \fn{batch\_main.log}.

    \begin{lstlisting}[language=bash]
python -m usseg.main
    \end{lstlisting}

  \subsection{Configuration}
    Configuration is read from \fn{config.toml} in the working directory. The
    keys are \fn{root\_dir}, the directory tree to search for scans;
    \fn{output\_dir}, where rendered PNG outputs are written; and a
    \texttt{[pickle]} table naming the intermediate serialisation files. Of the
    three pickle keys, \fn{likely\_us\_images} and \fn{segmented\_data} are used;
    \fn{patient\_paths} is declared but never read by any module.\\
    \\
    The file is loaded by relative path, so all commands must be run from the
    repository root.

  \subsection{Environment prerequisites}
    Text extraction depends on Tesseract OCR, configured by
    \fn{setup\_tesseract} in \fn{setup\_environment.py}. On Windows the
    executable is expected at
    \texttt{C:/Program Files/Tesseract-OCR/tesseract.exe}. Because every metric
    read from the image panel and every axis label depends on OCR, an
    incorrectly configured Tesseract installation causes the text and axis
    calibration stages to fail even though segmentation may still succeed.

\newpage
\section{Input Classification and Modality Dispatch}
\label{sec:dispatch}
  \subsection{Extension-based dispatch}
    Both entry points classify each input by file extension and set two mutually
    exclusive flags, \fn{us\_image} and \fn{us\_dicom}. Raster extensions are
    \texttt{.jpg}, \texttt{.jpeg}, \texttt{.png}, \texttt{.bmp}, \texttt{.gif},
    \texttt{.tif}, \texttt{.tiff} and \texttt{.webp}; DICOM extensions are
    \texttt{.dcm} and \texttt{.dicom}. Any other extension raises
    \fn{ValueError} in the single-image path, and is skipped with a warning in
    the batch path.\\

    Raster inputs are loaded twice, into a Pillow image and an OpenCV BGR array.
    The Pillow copy is converted to RGB and used for the colour-based operations
    (text colour isolation and coarse segmentation); the OpenCV array is used for
    thresholding, contour work, OCR input and all three tracing methods. DICOM
    inputs are read with \fn{extract\_dicom\_metadata} and
    \fn{extract\_doppler\_from\_dicom}, which return the calibration metadata and
    an image pair cropped to the Doppler region.

  \subsection{Why the two branches diverge}
    The single most important structural difference is the source of
    calibration. Table~\ref{tab:modality} summarises how the stages differ.

    \begin{table}[h]
    \centering
    \small
    \begin{tabular}{@{}lll@{}}
    \toprule
    Stage & Raster image & DICOM \\
    \midrule
    Metric text        & OCR of the on-screen panel & Not attempted \\
    Region bounds      & \fn{initial\_segmentation}  & Region metadata tags \\
    Velocity scale     & Detected axis ticks/labels  & \fn{PhysicalDeltaY} \\
    Time scale         & Inferred from heart rate    & \fn{PhysicalDeltaX} \\
    Zero line          & \fn{estimate\_zero\_line\_y} & \fn{ReferencePixelY0} \\
    Yellow exclusion   & Enabled                     & Disabled \\
    Output metrics     & OCR compared with digitised & Digitised only \\
    \bottomrule
    \end{tabular}
    \caption{Stage-by-stage differences between the two processing branches.}
    \label{tab:modality}
    \end{table}

    Everything from Section~\ref{sec:segmentation} onwards, that is the refined
    mask and all three tracing methods, is shared between the branches. The
    branches differ before segmentation, in how the region of interest and the
    calibration are established, and after segmentation, in how pixel
    coordinates are converted into physical units and which metrics are
    reported.

\newpage
\section{Stage 1: Text Extraction and OCR Repair}
\label{sec:ocr}
  \emph{Raster images only.} This stage is skipped entirely for DICOM input,
  where the metric panel is generally absent from the extracted Doppler region.\\
  \\
  The purpose of this stage is to recover the clinical measurements that the
  ultrasound machine has already printed on the screen. These values serve two
  roles: they are reported as the primary metric table, and they act as the
  reference against which the digitised waveform is later scored, which is what
  ultimately selects one of the three segmentation models.

  \subsection{Metric panel isolation: \fn{colour\_extract\_vectorized}}
    On the supported machines the metric panel is rendered in a distinctive
    yellow. Rather than thresholding brightness, the code isolates that colour
    geometrically: the target colour defines the axis of a cylinder in RGB space,
    and a pixel is retained if it falls inside that cylinder.\\
    \\
    The target is $[255, 255, 100]$ with both \fn{cyl\_length} and
    \fn{cyl\_radius} set to $95$. The axis direction is obtained from
    \fn{append\_spherical\_1point}, which expresses the target colour in
    spherical coordinates, and the cylinder is centred on the target with
    endpoints \fn{start} and \fn{end}. A pixel $p$ is inside when it lies between
    the two end planes and within the radius of the axis:

    \[
      (p-\fn{start})\cdot v \geq 0, \qquad
      (p-\fn{end})\cdot v \leq 0, \qquad
      \|(p-\fn{start}) \times v\| \leq \fn{cyl\_radius}\,\|v\|,
    \]

    where $v = \fn{end} - \fn{start}$. Matching pixels are set to white and all
    others to black, producing a binary mask image. Because the test is a
    distance in colour space rather than a per-channel range, it tolerates the
    brightness variation of anti-aliased text.\\
    \\
    An earlier non-vectorised implementation, \fn{colour\_extract}, remains in the
    module but is no longer used.

  \subsection{OCR and line assembly: \fn{text\_from\_greyscale}}
    This function performs OCR on the isolated panel and turns the result into a
    dataframe with columns \texttt{Line}, \texttt{Word}, \texttt{Value} and
    \texttt{Unit}. It also returns a \fn{Fail} flag.

    \paragraph{Spatial pre-masking.} Before OCR, the mask is cropped by blanking
    the left two-thirds of the width and the bottom 55\,\% of the height, because
    on the supported layout the metric panel occupies the upper right of the
    screen. This substantially reduces spurious detections from the waveform area.

    \paragraph{Primary OCR pass.} Tesseract is invoked through
    \fn{image\_to\_data} rather than \fn{image\_to\_string}, so that per-word
    bounding boxes and confidences are available for the line-grouping step:

    \begin{lstlisting}
--oem 1 --psm 3 -c tessedit_char_blacklist=l,!_|=$
    \end{lstlisting}

    \fn{oem 1} selects the LSTM engine and \fn{psm 3} fully automatic page
    segmentation. The character blacklist suppresses glyphs that are commonly
    confused with digits or that never appear in a metric label.

    \paragraph{Line grouping.} Words must be reassembled into lines before a
    label can be matched. The vertical centre of each word box is computed and
    clustered with DBSCAN over that single coordinate, using
    $\fn{eps} = 5$ pixels and $\fn{min\_samples} = 2$. Words within a cluster are
    sorted by horizontal position and joined with spaces. A consequence of
    $\fn{min\_samples} = 2$ is that an isolated single-word line forms no cluster
    and is therefore dropped from this source. To compensate, a second pass with
    \fn{image\_to\_string} provides newline-delimited lines, and the two sources
    are combined into a deduplicated union.

    \paragraph{The \fn{Fail} flag.} \fn{Fail} is set to $1$ when Tesseract
    returns fewer than 30 text tokens, which is treated as evidence that the
    image is not a metric-panel scan. It is a heuristic on text quantity, not a
    per-metric validity flag. Both entry points log a warning and continue when
    \fn{Fail} is set, but skip \fn{plot\_correction}, so a scan flagged this way
    still yields a segmentation and a trace but no OCR comparison and hence no
    informed model selection.

  \subsection{Target label sets and vessel families}
    Matching is constrained to a vessel family so that, for example, an umbilical
    label cannot be matched on a uterine scan. The recognised families are held in
    \fn{METRIC\_FAMILY\_PREFIXES}, in order:

    \begin{lstlisting}
("Lt Ophthalmic", "Rt Ophthalmic", "Lt MCA", "Rt MCA",
 "Lt Ut", "Rt Ut", "Umb", "DV")
    \end{lstlisting}

    The order matters because the more specific prefixes are tested first.
    \fn{\_detect\_metric\_family\_from\_lines} counts how many OCR lines match
    each prefix and selects the highest-scoring family; ophthalmic lines
    additionally increment the count through \fn{\_line\_looks\_ophthalmic}. If no
    family is detected the full target list is used and a warning is logged.

    \begin{table}[h]
    \centering
    \small
    \begin{tabular}{@{}lll@{}}
    \toprule
    Vessel & Prefix & Metrics sought \\
    \midrule
    Uterine (left/right) & \texttt{Lt Ut-}, \texttt{Rt Ut-} & PS, ED, S/D, PI, RI, MD, TAmax, HR \\
    Umbilical            & \texttt{Umb-}   & as above \\
    Middle cerebral      & \texttt{Lt MCA-}, \texttt{Rt MCA-} & as above \\
    Ductus venosus       & \texttt{DV-}    & S, D, a, TAmax, S/a, a/S, PI, PLI, PVIV, HR \\
    Ophthalmic           & \texttt{Lt/Rt Ophthalmic A.} & PS ($\times 2$), ED, PI, RI, PS/ED, ED/PS \\
    \bottomrule
    \end{tabular}
    \caption{Vessel families and the metrics sought for each.}
    \label{tab:families}
    \end{table}

    Two label lists exist per family: a plain list and an
    \emph{extended} list carrying unit suffixes, for example
    \texttt{"Lt Ut-PS cm/s"}. The extended list is used only by the final fuzzy
    stage described below.

  \subsection{Label matching}
    Non-ophthalmic families are matched by a four-stage cascade, each stage
    operating only on the targets still unmatched by the previous one. This
    ordering is deliberate: cheap exact tests run first, and progressively looser
    tests are applied only where necessary.

    \begin{enumerate}
      \item \textbf{Exact substring.} The target must appear literally in the
        line.
      \item \textbf{Subsequence.} The target's characters, ignoring spaces and
        hyphens, must appear in order within the line. This absorbs dropped
        characters without permitting reordering.
      \item \textbf{Prefix Levenshtein.} The first \texttt{len(word)} characters
        of the line are compared with the target using
        \fn{rapidfuzz.distance.Levenshtein}, and accepted at an edit distance of
        $7$ or less.
      \item \textbf{Extended-label distance matrix.} Digits are stripped from the
        line and the edit distance to each extended label is computed. If the
        expected metric suffix (for example \texttt{PS} or \texttt{PI}) appears in
        the line, the distance is biased by $-2$. Each remaining target then takes
        the line with the smallest non-zero distance, assigned one-to-one so that
        a line cannot serve two targets.
    \end{enumerate}

    \paragraph{Ophthalmic matching.} Ophthalmic scans bypass the cascade entirely
    and use \fn{\_match\_ophthalmic\_lines}. The reason is structural: ophthalmic
    panels report two peak-systolic measurements, so labels cannot be identified
    by text alone and position must be used. Each line is classified by
    \fn{\_ophthalmic\_metric\_kind} into \texttt{PS}, \texttt{ED}, \texttt{PI},
    \texttt{RI}, \texttt{PS/ED} or \texttt{ED/PS}, and \texttt{PS} lines are
    assigned to the numbered targets in the order encountered. Orphan numeric
    lines with no label are attached to the first PS row still holding zero.
    Laterality is inferred by \fn{\_detect\_ophthalmic\_side\_from\_ocr}, which
    counts standalone \texttt{Lt} and \texttt{Rt} tokens and resolves ties in
    favour of \texttt{Rt}. \textbf{This is not a comprehensive solution to Ophthalmic label matching and improvements could be made.}

  \subsection{Value parsing and numeric repair}
    Values are taken from the end of the matched line with an anchored regular
    expression that tolerates OCR-split decimals such as \texttt{"27 . 35"}, and
    captures any trailing non-numeric group as the unit. Spaces are removed from
    the captured number before conversion. A matched label with no parseable
    number yields a row with a zero value rather than being dropped, so that the
    metric's absence is visible downstream.\\
    \\
    There is no character-level substitution table (for instance mapping a
    misread \texttt{S} to \texttt{5}). Instead, misplaced decimal points are
    repaired by magnitude: a value that is implausible by an order of magnitude is
    divided down into its physiological range. \fn{check\_pi\_value} divides a
    pulsatility index by 10 in the range $3$--$10$ and by 100 in the range
    $10$--$200$; the PS/ED normalisation applies the same principle for
    velocities, dividing by 10 between $250$ and $1000$ and by 100 between $1000$
    and $10000$; \fn{check\_S\_D\_value} does likewise for ductus venosus values.
    The approach is effective because the true ranges of these quantities are
    narrow and well separated from their decimal-shifted counterparts.

  \subsection{Heart rate refinement}
    Heart rate deserves special treatment because it is the only quantity that
    calibrates the raster time axis (Section~\ref{sec:digitisation}); an error
    here propagates into every time-dependent result.\\
    \\
    After the first pass, \fn{refine\_hr\_from\_local\_roi} re-reads the heart
    rate from a locally enhanced crop. The ROI around the HR line is padded by
    $\max(15,\,35\%)$ of its width and $\max(10,\,80\%)$ of its height, then
    equalised, blurred, Otsu-thresholded and upscaled by a factor of two before a
    constrained OCR pass using \fn{psm 7} with a character whitelist, falling back
    to \fn{psm 6}. Candidates are accepted only within $20$--$240$\,bpm.\\
    \\
    The reconciliation rule is conservative. The refined value replaces the first
    pass only if the first-pass value is implausible; if both are plausible but
    differ by more than $3$\,bpm the first-pass value is kept and a warning is
    logged. Finally the absolute value is taken, since a spurious minus sign is a
    common OCR artefact.

  \subsection{Metric consistency checks}
    The final step exploits the algebraic relationships between the reported
    metrics to detect and repair OCR errors. Three variants exist, dispatched by
    family: \fn{metric\_check} for uterine, umbilical and MCA scans,
    \fn{metric\_check\_ophthalmic} for ophthalmic, and \fn{metric\_check\_dv} for
    ductus venosus. All of them preserve the original readings in a
    \texttt{Raw Value} column before making corrections, so that repairs remain
    auditable.\\
    \\
    \fn{metric\_check} recomputes the three derived quantities from the extracted
    peak-systolic and end-diastolic velocities,
    \[
      (S/D)_{\text{calc}} = \frac{PS}{ED}, \qquad
      RI_{\text{calc}} = \frac{PS-ED}{PS}, \qquad
      TA_{\text{max,calc}} \approx \frac{PS + 2\,ED}{3},
    \]
    and compares each with the extracted value at a tolerance of $0.5$. Time
    averaged maximum velocity is treated specially: it is accepted if its
    magnitude simply lies between $ED$ and $PS$, because the closed-form estimate
    above is only an approximation.\\
    \\
    The outcome depends on how many of the three agree. If all three agree,
    nothing is changed. If one or two agree, the disagreeing or missing fields are
    recomputed from the corrected velocities. If none agree, the code assumes $PS$
    is correct and back-solves $ED$ from the reported ratio or resistive index;
    should both resulting candidates fall below $2$, it reverses the assumption,
    treats $ED$ as correct and recomputes $PS$. This is the most aggressive repair
    in the pipeline and the primary reason the raw values are retained.\\
    \\
    The ophthalmic and ductus venosus variants follow the same pattern with
    different tolerances and formulae: ophthalmic uses $0.5$ for the ratio and
    $0.1$ for the resistive index, while ductus venosus uses $0.2$ for $S/a$ and
    $PI$ and a looser $2.0$ for the approximated time-averaged velocity. The
    ductus venosus path additionally applies pre-swap heuristics in
    \fn{text\_from\_greyscale} to correct values that OCR has attributed to the
    wrong label.

\newpage
\section{Stage 2: Coarse Localisation of the Doppler Region}
\label{sec:coarse}
  \subsection{Raster: \fn{initial\_segmentation}}
    \fn{initial\_segmentation} finds the approximate extent of the Doppler
    waveform panel in a raster screenshot. It exploits the fact that the
    waveform and its surrounding trace are rendered in near-neutral greys, so
    they can be separated from coloured annotation and from the dark background
    by a colour-neutrality test rather than by brightness alone.\\
    \\
    The function iterates over every pixel of the RGB image and retains the
    pixel only if the spread across the three channels is small and the pixel is
    reasonably bright:

    \[
      \left(\max(r,g,b) - \min(r,g,b)\right) < 100
      \quad\text{and}\quad
      \max(r,g,b) > 120 .
    \]
    Pixels satisfying both conditions are set to white and all others to black.
    A band at the top of the image is then forced to black to suppress the
    header region: rows above 500 for images taller than 600 pixels, and rows
    above 20 otherwise.\\
    \\
    The resulting binary image is cleaned by removing small objects and small
    holes at a size threshold of 199 pixels, followed by two erosions and one
    dilation. The erosions detach tick marks that touch the waveform body and
    the dilation recovers part of the resulting loss. Contours are then extracted
    with \fn{skimage.measure.find\_contours}, and the returned bounds
    \fn{Xmin}, \fn{Xmax}, \fn{Ymin} and \fn{Ymax} are the extremes across
    \emph{all} contours, that is the bounding box of every surviving object
    rather than of a single selected component.\\
    \\
    Two practical consequences follow. First, because the bounding box spans all
    surviving contours, residual bright artefacts inside the search area widen
    the reported region. Second, the per-pixel Python loop makes this one of the
    slowest steps in the pipeline; it is a natural target for vectorisation.

  \subsection{Raster: \fn{define\_end\_rois}}
    \fn{define\_end\_rois} derives the two side regions in which axis ticks and
    labels are later sought. Given the coarse bounds and the image dimensions it
    returns two \texttt{[Xmin, Xmax, Ymin, Ymax]} lists:

    \begin{description}
      \item[Left ROI] spans horizontally from column $0$ to $\fn{Xmin}-1$, that
        is everything to the left of the waveform.
      \item[Right ROI] spans horizontally from \fn{Xmax} to the right-hand edge
        of the image.
    \end{description}
    Both regions share the same vertical extent. The top is set to
    $\fn{Ymin}-75$ when $\fn{Ymin}-125 > 0$, and to row $1$ otherwise; the
    bottom is the full image height. The upward extension is deliberate: axis
    labels are frequently printed slightly above the top of the traced waveform,
    so an ROI clipped exactly at \fn{Ymin} would omit the highest label.

  \subsection{DICOM: region bounds from metadata}
    The DICOM branch does not call either function above. The waveform bounds are
    taken directly from the ultrasound region metadata as
    \fn{RegionLocationMinX0}, \fn{RegionLocationMaxX1},
    \fn{RegionLocationMinY0} and \fn{RegionLocationMaxY1}, and the zero-velocity
    row is computed as $\fn{RegionLocationMinY0} + \fn{ReferencePixelY0}$. The
    bounds are therefore exact rather than estimated, which removes a whole
    class of failure that the raster branch must tolerate.

\newpage
\section{Stage 3: Axis Calibration}
\label{sec:axes}
  Calibration establishes the mapping from image row to physical velocity. This
  is the stage where the two modalities differ most sharply: for DICOM it is a
  metadata lookup, whereas for raster images the entire scale must be recovered by
  locating tick marks and reading their printed values by OCR.

  \subsection{Raster: tick detection with \fn{search\_for\_ticks}}
    \fn{search\_for\_ticks} is called once per side with \fn{side} set to
    \texttt{"Left"} or \texttt{"Right"}, and returns twelve values. The important
    ones are \fn{Cs} (tick contours), \fn{CenPoints} (tick centres in ROI-local
    coordinates), \fn{ROIAX} (the axis ROI cropped so that it contains the labels
    but not the tick strokes) and \fn{TYLshift} (the offset used for that crop).
    The remainder are intermediate images and pass-through arguments consumed by
    \fn{search\_for\_labels}.\\
    \\
    The algorithm is:

    \begin{enumerate}
      \item Crop the side ROI produced by \fn{define\_end\_rois} and threshold the
        greyscale image at an intensity of 127.
      \item Find contours in the ROI and draw them into a working image.
      \item \textbf{Right side only:} re-label the ROI and retain only components
        with an area between 5 and 200 pixels and a height below 5 pixels, that is
        objects shaped like short horizontal strokes. Columns are then blanked if
        they intersect more than 20 such objects or none at all. This suppresses
        the dense text of the adjacent panel, which is a problem on the right side
        but not the left.
      \item Re-find contours on the filtered image and compute each contour's
        centre.
      \item \textbf{Select the tick column.} For every column in the ROI, count
        the contours whose $x$ extent includes that column, giving a
        one-dimensional profile. \fn{scipy.signal.find\_peaks} is applied to that
        profile with $\fn{height} = 3$, requiring at least three aligned objects
        for a column to be considered.
      \item Retain the contours lying on the selected column, recording their
        centres and endpoints.
      \item \textbf{Left side only:} reject the worst endpoint outlier using a
        modified $z$-score against the median with $m = 8.0$.
      \item Compute \fn{TYLshift} and re-crop \fn{ROIAX} so that OCR sees the
        numeric labels without the tick strokes, which would otherwise be read as
        hyphens.
    \end{enumerate}
    %TODO: Add a figure showing the tick detection process.
    %TODO: Explain terms like TYLshift, ROIAX, etc.
    \paragraph{Asymmetry in column selection.} The two sides choose the tick
    column by different criteria, and this is the most consequential difference
    between them. For the right axis the strongest peak is chosen by the product
    of peak height and peak width:

    \begin{lstlisting}[language=Python]
peak_wid = signal.peak_widths(all, peaks)
maxID = np.argmax(vals["peak_heights"] * peak_wid[0])
    \end{lstlisting}

    For the left axis the selection is instead:

    \begin{lstlisting}[language=Python]
maxID = np.argmax(peaks)
    \end{lstlisting}

    Because \fn{peaks} contains column indices in ascending order, taking its
    \fn{argmax} selects the \emph{last} entry, that is the right-most candidate
    column in the left ROI. This is the column nearest the waveform, which is
    where left-axis ticks are expected, so the behaviour is reasonable; but it is
    expressed as a maximum over positions rather than as an explicit choice of the
    last peak, which obscures the intent and makes the line easy to misread.

    \paragraph{Failure when no ticks exist.} Both selections index \fn{peaks}
    unconditionally. On a scan with no ticks on a given side the profile has no
    peaks, \fn{peaks} is empty, and the call raises
    \texttt{ValueError: attempt to get argmax of an empty sequence}. This is the
    normal, expected outcome for scans that print labels on one side only, and it
    is caught by the caller; see Section~\ref{sec:axistolerance}.

  \subsection{Raster: label detection and pairing with \fn{search\_for\_labels}}
    This function reads the numeric tick values and associates each with a tick
    position. It returns the cropped ROI, the list of label strings
    (\fn{number}), their paired positions in full-image coordinates
    (\fn{positions}), and a full-size debug canvas.

    \paragraph{Threshold sweep with a plausibility stop.} Rather than committing
    to one binarisation, the function sweeps the threshold from 100 to 185 in
    steps of 5, running OCR at each value with

    \begin{lstlisting}
--psm 11 -c tessedit_char_whitelist=-0123456789
    \end{lstlisting}
    and stops as soon as every parsed number is divisible by 5. This exploits a
    strong prior about Doppler displays: velocity axes are labelled in multiples
    of five, so divisibility is an inexpensive proxy for a clean read. The
    whitelist restricts output to digits and the minus sign, and \fn{psm 11}
    treats the crop as sparse text.

    \paragraph{Box handling.} Tesseract emits repeated rows per box, which are
    deduplicated by an index test; the first box is skipped, and empty or
    hyphen-only labels are discarded. Label centres are computed from the bounding
    boxes, with \fn{TYLshift} added back on the right side to undo the earlier
    crop. If any two label centres lie within 20 pixels of each other, the code
    logs \texttt{"Characters too close"}, indicating that two labels may have been
    merged.

    \paragraph{Pairing.} Each label centre is matched to its nearest tick centre
    by Euclidean distance using \fn{scipy.spatial.distance.cdist}, and the
    matched tick position is offset back into full-image coordinates. Finally
    \fn{correct\_number\_format} strips a trailing hyphen from values such as
    \texttt{"40-"}, an artefact of a tick stroke that survived the crop.

    \paragraph{Fallback.} If pairing fails, the function re-runs OCR and uses the
    label box centres themselves as the tick positions. This is less accurate,
    because a label's centre is horizontally offset from its tick, but it keeps
    calibration possible when tick detection has partially failed. The
    \fn{cdist} call is also the origin of the
    \texttt{ValueError: XA must be a 2-dimensional array} seen on scans with no
    labels on one side, since an empty list of centre points has no second
    dimension.

  \subsection{Raster: validation of ticks}
    \fn{validate\_axis\_ticks} cleans a single side's ticks. It never mutates its
    inputs and returns a corrected copy. Values that cannot be parsed as floats
    are skipped, and if fewer than two ticks remain, or the value and position
    counts disagree, the inputs are returned unchanged.\\
    \\
    Ticks are sorted by pixel row and then repeatedly filtered until stable:

    \begin{itemize}
      \item Ticks sharing an identical row are collapsed: if they carry the same
        value the last is kept, and if they conflict the whole cluster is dropped.
      \item The slope $\Delta \text{value} / \Delta y$ is computed between
        successive ticks. Slopes are clustered within a relative tolerance
        \fn{spacing\_tol} $= 0.2$, and the largest cluster is taken as dominant.
      \item A single anomalous slope at either end drops the corresponding end
        tick; exactly two consecutive anomalous slopes drop the tick between them;
        a specific three-slope pattern drops both an interior tick and an
        endpoint. Any more complex pattern only logs a warning and stops, on the
        principle that it is safer to keep suspect ticks than to guess.
    \end{itemize}
    One caveat worth recording: the function computes monotonicity flags for the
    ordered values but does not act on them. A non-monotonic set of tick values is
    therefore not rejected on that basis alone, and is only corrected if it also
    produces an anomalous slope pattern.

  \subsection{Raster: cross-checking the two axes}
    \fn{validate\_axis\_pair} compares the sides where both were detected. It
    builds a value-to-row map for each side, intersects the tick values common to
    both, and warns if the mean row for a shared value differs by more than
    \fn{y\_tol\_pixels} $= 5.0$. When the sides share no values the result is
    inconclusive and reported as agreement. The function is purely diagnostic: it
    returns a boolean and changes no data, and its result is used only to decide
    which side to trust during digitisation.

  \subsection{Raster: zero-line estimation}
    \fn{estimate\_zero\_line\_y\_axis} finds the row corresponding to zero
    velocity for one side. If any tick value is zero within $10^{-3}$ it returns
    the mean row of those ticks. Otherwise, given at least two ticks, it fits a
    straight line of row against value with \fn{numpy.polyfit} and returns the
    intercept at zero. The fit matters because it recovers the baseline even when
    it is off-screen or unlabelled, which is common when a display shows only
    positive velocities.\\
    \\
    \fn{estimate\_zero\_line\_y} then fuses the two sides. With one side available
    it returns that estimate. With both, it prefers whichever side has at least
    three ticks when the other has fewer; if the sides are comparable it warns
    when they differ by more than 5 pixels and returns their mean. The result is
    passed to \fn{segment\_refinement}, where it decides waveform orientation
    (Section~\ref{sec:keep}). A \fn{None} result is tolerated and triggers the
    centroid-based fallback described there.

  \subsection{DICOM: calibration from metadata}
    The DICOM branch performs none of the above. \fn{extract\_dicom\_metadata}
    reads the Sequence of Ultrasound Regions, tag $(0018,6011)$, and selects the
    first region whose Region Spatial Format and Region Data Type are both $3$,
    identifying the pulsed-wave spectral region. From that region it reads the
    bounding box, the reference pixel, the reference physical values and the
    physical deltas, listed in Table~\ref{tab:dicomtags}.\\
    \\
    \begin{table}[h]
    \centering
    \small
    \begin{tabular}{@{}lll@{}}
    \toprule
    Key & Tag & Role \\
    \midrule
    \fn{RegionLocationMinX0}          & $(0018,6018)$ & ROI left \\
    \fn{RegionLocationMinY0}          & $(0018,601A)$ & ROI top \\
    \fn{RegionLocationMaxX1}          & $(0018,601C)$ & ROI right \\
    \fn{RegionLocationMaxY1}          & $(0018,601E)$ & ROI bottom \\
    \fn{ReferencePixelX0}             & $(0018,6020)$ & Origin column \\
    \fn{ReferencePixelY0}             & $(0018,6022)$ & Origin row (baseline) \\
    \fn{ReferencePixelPhysicalValueX} & $(0018,6028)$ & Physical value at origin \\
    \fn{ReferencePixelPhysicalValueY} & $(0018,602A)$ & Physical value at baseline \\
    \fn{PhysicalDeltaX}               & $(0018,602C)$ & Units per column \\
    \fn{PhysicalDeltaY}               & $(0018,602E)$ & Units per row \\
    \bottomrule
    \end{tabular}
    \caption{DICOM tags used for calibration.}
    \label{tab:dicomtags}
    \end{table}
    Only the header is read, using \fn{stop\_before\_pixels=True}. Missing
    elements resolve to \fn{None} rather than raising, and any exception causes
    the function to return an empty dictionary; in that case the absence of
    bounding-box keys will surface later as an arithmetic error during
    digitisation. The physical units tags $(0018,6024)$ and $(0018,6026)$ are
    recorded but not consulted during conversion, so the code implicitly assumes
    the conventional units of centimetres per second and seconds.

  \subsection{Tolerance of missing axes}
  \label{sec:axistolerance}
    Both entry points wrap each side's tick and label search in its own
    \fn{try}/\fn{except}, initialising \fn{Lnumber}, \fn{Lpositions},
    \fn{Rnumber} and \fn{Rpositions} to \fn{None} beforehand. A side that fails
    leaves its variables as \fn{None} and the pipeline proceeds on the other
    side. This is what allows scans that label only one axis to be processed at
    all, and it is by design rather than by accident.\\
    \\
    The cost is diagnostic noise. In \fn{segment\_files.py} the handlers call
    \fn{traceback.print\_exc()} in addition to logging, so an entirely expected
    single-sided scan still prints full tracebacks for every image processed. The
    exceptions are handled, but they are not silenced, which makes a healthy batch
    run look as though it is failing.

\newpage
\section{Stage 4: Waveform Segmentation and the Three Tracing Methods}
\label{sec:segmentation}
  This is the algorithmic core of the package and it is shared by both
  modalities. It proceeds in two parts. First a single refined binary mask of the
  waveform body is produced. Then \emph{three independent methods} extract a
  one-pixel-per-column trace of the waveform envelope. All three are computed
  for every scan; the choice between them is deferred to
  Section~\ref{sec:selection}.

  \subsection{Orchestration: \fn{segment\_refinement}}
    \fn{segment\_refinement} is the single entry point for this stage and is
    called identically by both entry points, except for the
    \fn{grow\_forbid\_yellow} flag which is \texttt{True} for raster images and
    \texttt{False} for DICOM. Its sequence is:

    \begin{enumerate}
      \item Build the refined binary mask with
        \fn{refine\_waveform\_segmentation}.
      \item Convert the image to greyscale and construct the permitted region
        \fn{allowed} with \fn{allowed\_mask\_from\_roi\_bounds}.
      \item If \fn{grow\_forbid\_yellow} is set, build a forbidden-region mask of
        instrument yellow with \fn{hsv\_yellow\_tick\_mask\_bgr}.
      \item Shrink the refined mask into a conservative seed with
        \fn{\_shrink\_seed\_for\_region\_grow}.
      \item Grow that seed with \fn{constrained\_region\_grow}.
      \item Derive all three envelope traces with \fn{compute\_top\_curve}.
    \end{enumerate}
    It returns seven values: the refined mask, and a mask/coordinate pair for
    each of the morphological, ray-traced and grown traces. The latter two pairs
    are \fn{None} if their method failed or produced an empty result, so all
    downstream code must tolerate missing traces.

  \subsection{The shared refined mask: \fn{refine\_waveform\_segmentation}}
    This function produces the binary waveform body that both the morphological
    method and the region-growing seed are derived from. Unlike
    \fn{initial\_segmentation} it works on the greyscale OpenCV image and uses a
    fixed intensity threshold:

    \begin{enumerate}
      \item Convert BGR to greyscale and threshold at an intensity of 30.
      \item Zero everything outside the region of interest. Note the asymmetric
        vertical margin: the top boundary is $\fn{Ymin}-50$ rather than
        \fn{Ymin}, deliberately admitting rows above the coarse bound so that
        systolic peaks clipped by Stage~2 are recoverable.
      \item Remove small objects and holes (199 pixels), erode twice, dilate.
      \item Dilate again, remove holes up to 999 pixels, apply a morphological
        closing, then a $3\times3$ median filter.
      \item Label connected components and discard everything smaller than the
        second-largest area plus 10, which in practice retains only the main
        waveform body.
      \item Blur with a Gaussian of $\sigma = 7$ and re-threshold at $0.5$,
        smoothing the outline so that the envelope extracted from it is not
        dominated by pixel-level noise.
    \end{enumerate}

  \subsection{Envelope orientation: upper or lower}
  \label{sec:keep}
    Doppler waveforms may be displayed above or below the baseline, so before any
    envelope is extracted the code must decide which side of the waveform body is
    the diagnostic envelope. This decision, held in the variable \fn{keep} with
    value \texttt{"upper"} or \texttt{"lower"}, is made once in
    \fn{\_morphological\_top\_curve\_from\_mask} and then \emph{reused} for the
    ray tracer, ensuring that all methods trace the same side of the waveform.\\
    \\
    When a zero line is available the mask outline rows are counted above and
    below it, and the waveform is treated as inverted when more of the outline
    lies below the baseline. Rows on the non-diagnostic side are then cleared.
    When no zero line is available the code falls back to the centroid row of the
    first labelled region and validates the choice with
    \fn{check\_inverted\_curve}, which declares the trace inverted if the
    retained outline spans less than a fraction \texttt{tol} $= 0.25$ of the
    waveform height; if so, the opposite side is kept instead.

  \subsection{Method 1: Morphological envelope}
  \label{sec:morph}
    Implemented by \fn{\_morphological\_top\_curve\_from\_mask}. This is the
    original method and the only one guaranteed to be produced.\\
    \\
    The mask is eroded once and the erosion subtracted from the original,
    $\text{outline} = \text{mask} - \text{erosion}(\text{mask})$, leaving a
    one-pixel shell around the waveform body. The non-diagnostic side is cleared
    as described in Section~\ref{sec:keep}, and
    \fn{keep\_one\_pixel\_per\_column} reduces the shell to a single pixel per
    column, taking the visually uppermost or lowermost pixel according to
    \fn{keep}.\\
    \\
    The method is fast, deterministic and entirely dependent on the quality of
    the refined mask. Because the mask has been Gaussian-smoothed, the resulting
    trace is smooth but can be blunted at sharp systolic peaks, and any
    thresholding error in the mask is inherited directly.

  \subsection{Method 2: Ray tracing}
  \label{sec:ray}
    Implemented by \fn{ray\_trace\_waveform\_segmentation} and used as the
    primary metric source downstream. Rather than deriving the envelope from a
    binary mask, this method re-examines the original image column by column,
    which allows it to follow a faint envelope that global thresholding would
    lose.\\
    \\
    Its region of interest is computed by
    \fn{ray\_trace\_roi\_from\_refined\_mask} from the bounding box of the
    refined mask, padded by $\max(8, 3\%)$ of the box width horizontally,
    $\max(18, 35\%)$ of its height above and $\max(6, 6\%)$ below, then inset by
    15 pixels horizontally. The generous top padding exists so that peaks
    excluded from the refined mask remain inside the search area. A further hard
    inset of 20 pixels per side is applied inside the tracer itself to exclude
    border artefacts.\\
    \\
    The algorithm then proceeds as follows.

    \begin{enumerate}
      \item \textbf{Adaptive signal/background separation.} All ROI pixels are
        clustered in BGR colour space with $k$-means, $k=3$. The cluster with the
        lowest mean greyscale is declared background and the remaining two
        clusters become signal candidates. Because the split is learned per
        image, this adapts to scans of differing brightness instead of relying on
        a fixed threshold.
      \item \textbf{Instrument-yellow removal.} Pixels identified by
        \fn{hsv\_yellow\_tick\_mask\_bgr} (hue $18$--$40$, saturation $>90$,
        value $>90$, followed by an opening with a disk of radius 1) are removed
        from the signal, so that the machine's own yellow peak markers are not
        mistaken for the envelope.
      \item \textbf{Continuity-constrained column scan.} Columns are traversed
        left to right. Each column is first searched only within
        $\pm\,\fn{max\_col\_step\_y}$ rows of the previous column's pick, which
        enforces vertical continuity and prevents the trace jumping to an
        unrelated structure. Within that window,
        \fn{\_pick\_y\_from\_column\_signal} prefers the first run of three
        consecutive signal pixels, taking the uppermost or lowermost such run
        according to \fn{keep}; if no run exists it falls back to the brightest
        pixel in the window. If the banded search fails entirely the whole column
        is searched and the result is clipped back into the permitted step.
      \item \textbf{Gap filling and smoothing.} Columns with no pick are filled
        by linear interpolation, isolated outliers are suppressed by a Hampel
        filter (half-window 2, $3\sigma$), and a width-3 median filter is
        applied. The narrow kernels are a deliberate compromise: they remove
        spikes while preserving the steep genuine slope into the diastolic foot.
      \item The trace is rasterised, dilated by a disk of radius 1 and mapped
        back into full-image coordinates.
    \end{enumerate}
    The step limit \fn{max\_col\_step\_y} defaults to
    $\max(30, \text{ROI height}/25)$ and is clamped to at least 3. It is the
    method's most consequential parameter: too small and the trace cannot follow
    a steep systolic upstroke, too large and continuity no longer constrains it.

  \subsection{Method 3: Constrained region growing}
  \label{sec:grow}
    Implemented by \fn{constrained\_region\_grow}. This method exists to recover
    parts of the waveform that the fixed threshold in
    \fn{refine\_waveform\_segmentation} misses, most notably the dimmer signal in
    the diastolic troughs, by expanding outward from a confident core rather than
    classifying each pixel independently.

    \paragraph{Seed preparation.} The refined mask is first shrunk by
    \fn{\_shrink\_seed\_for\_region\_grow}, which erodes it with a disk of radius
    \fn{GROW\_SEED\_EROSION\_RADIUS} $= 10$ and intersects the result with the
    permitted region. Shrinking matters because the seed statistics drive the
    acceptance test: an over-thick seed biases the intensity model towards the
    bright waveform core and prevents growth into the troughs. If erosion removes
    every seed pixel the function logs a warning and falls back to the
    unshrunken intersection.

    \paragraph{Constraints.} Growth is restricted to \fn{allowed}, the ROI mask
    produced by \fn{allowed\_mask\_from\_roi\_bounds} (which applies the same
    $\fn{Ymin}-50$ top margin as the refined mask). A separate
    \fn{forbidden\_mask} may additionally block expansion without excluding seed
    pixels; for raster images this carries the instrument-yellow mask, and for
    DICOM it is disabled via \fn{grow\_forbid\_yellow=False} because yellow is
    rare there and masking it removes real signal.

    \paragraph{Intensity model.} The greyscale image is normalised to $[0,1]$
    over the whole image. From the seed the mean and standard deviation are
    computed, and an absolute floor is derived as
    \[
      \fn{min\_intensity} = \max\!\left(
        \mu_{\text{seed}} - \max\!\left(5.0\,\sigma_{\text{seed}},\; 0.40\right),\;
        0.05 \right).
    \]
    The initial acceptance band is $\mu_{\text{seed}} \pm 0.43$, clamped by that
    floor and by 1.

    \paragraph{Growth.} A breadth-first flood fill runs from every seed pixel.
    The neighbourhood is a Chebyshev ball of radius
    \fn{GROW\_NEIGHBOUR\_CHEBYSHEV\_RADIUS} $= 2$, that is a $5\times5$
    neighbourhood, which lets the region bridge one-pixel gaps in a faint
    envelope; the \fn{connectivity} argument only applies when the radius is 1. A
    candidate is rejected outright below \fn{min\_intensity}. Otherwise, because
    \fn{update\_seed\_stats} defaults to \texttt{True}, the acceptance band is
    recomputed from the running mean and standard deviation of all accepted
    pixels as
    $\mu \pm \max(3.9\,\sigma, \fn{intensity\_tol})$, again clamped. This lets
    the region adapt as it moves into dimmer areas. Growth stops when the region
    reaches \fn{max\_area\_growth} $= 3.0$ times the seed area, which bounds both
    runtime and the extent of any leak into the background.

    \paragraph{Envelope extraction.} The grown mask is passed back through
    \fn{\_morphological\_top\_curve\_from\_mask}, so Method~3 differs from
    Method~1 only in the mask it operates on, not in how the envelope is thinned.

  \subsection{Dispatch and comparison: \fn{compute\_top\_curve}}
    \fn{compute\_top\_curve} produces all three traces and returns them as three
    mask/coordinate pairs. The morphological envelope is always computed first
    because it establishes the shared \fn{keep} orientation. The grown envelope is
    computed only if a non-empty grown mask was supplied, and the ray trace only
    if the original image and coarse $x$ bounds were supplied; the ray call is
    wrapped in a \fn{try}/\fn{except} so a tracing failure degrades to
    \fn{None} rather than aborting the scan.

    \begin{table}[h]
    \centering
    \small
    \begin{tabular}{@{}p{2.3cm}p{3.4cm}p{3.4cm}p{3.9cm}@{}}
    \toprule
     & Morphological & Ray tracing & Region growing \\
    \midrule
    Operates on & Refined binary mask & Original pixel data & Refined mask as seed \\
    Separation  & Fixed threshold of 30 & Per-image $k$-means, $k=3$ & Adaptive intensity band \\
    Envelope    & Erosion difference, thinned & Per-column pick & Erosion difference, thinned \\
    Continuity  & Implicit in the mask & Enforced by step limit & Implicit in connectivity \\
    Strength    & Robust, deterministic & Follows faint envelopes & Recovers dim troughs \\
    Weakness    & Inherits threshold error; blunts peaks & Sensitive to step limit & Can leak; area capped \\
    Always present & Yes & No & No \\
    \bottomrule
    \end{tabular}
    \caption{The three tracing methods compared.}
    \label{tab:methods}
    \end{table}

  \subsection{Debugging aids}
    Setting the module-level flag \fn{SHOW\_GROW\_DEBUG\_PLOTS}, or passing
    \fn{grow\_debug\_plots=True} to \fn{segment\_refinement}, produces
    diagnostic figures of the region-growing pipeline and of the traced overlays.
    \fn{ray\_trace\_waveform\_segmentation} has an equivalent
    \fn{debug\_plots} argument showing the clustered signal mask, the yellow
    exclusion, the accepted three-pixel runs and the final trace. These are off
    by default because the batch path saves a fixed set of figures instead.

\newpage
\section{Stage 5: Digitisation and Calibration}
\label{sec:digitisation}
  Digitisation converts a set of image coordinates into a physical waveform. The
  two modalities use different functions for this,
  \fn{plot\_digitized\_data\_single\_axis} for raster images and
  \fn{plot\_digitized\_data\_dicom} for DICOM, but both return the same
  seven-element tuple:

  \begin{lstlisting}[language=Python]
Xplot, Yplot, Ynought, Xplot_o, Yplot_o, Xplot_g, Yplot_g
  \end{lstlisting}
  where the unsuffixed pair is the primary series, \fn{\_o} the overlay and
  \fn{\_g} the region-grow series. \fn{Ynought} carries the zero-velocity
  reference. Keeping the signature identical is what allows the two branches to
  reconverge for the remaining stages.

  \subsection{Condensing the curve to one sample per column}
    Both branches begin with the same reduction. A trace arrives as a list of
    $(\text{row}, \text{col})$ pairs which may contain several rows for one
    column. The pairs are swapped to $(\text{col}, \text{row})$, grouped by
    column with \fn{pandas.groupby}, and the mean row per column is taken, after
    which the result is sorted by column. This guarantees that the digitised
    series is single-valued in time, which every downstream operation assumes.

  \subsection{Raster: choosing an axis and mapping rows to velocities}
    \fn{plot\_digitized\_data\_single\_axis} calibrates from a single vertical
    axis. A side is usable if it carries at least two ticks and its value and
    position counts agree. The choice is made as follows: if only one side is
    usable, it is used; if both are usable, \fn{validate\_axis\_pair} is consulted
    and the side with more ticks is selected either way, with a warning logged when
    the sides disagree. If neither side is usable the function returns empty
    series and logs an error. In other words tick count, not agreement, decides
    the outcome; agreement only affects logging. This is a deliberate choice to
    keep processing possible on partially labelled scans, and the disagreement
    warning is the signal to inspect such a scan manually.\\
    \\
    Calibration itself is a two-point fit rather than a least-squares fit over all
    ticks. The $(\text{row}, \text{value})$ pairs are sorted by row and the line is
    taken through the extreme pair only:

    \[
      a = \frac{v_{n-1} - v_0}{y_{n-1} - y_0}, \qquad
      b = v_0 - a\,y_0, \qquad
      \text{value} = a\,y + b .
    \]
    Using the endpoints maximises the lever arm and so minimises the effect of
    quantisation on the slope, but it also means a single mislocated endpoint tick
    corrupts the entire scale. This is precisely the failure that
    \fn{validate\_axis\_ticks} exists to prevent, which is why its endpoint-drop
    rules matter more than they might appear to.\\
    \\
    Interior ticks are used for validation but not for the fit; the code proceeds
    only if there are at least two ticks and the first and last rows differ.

  \subsection{Raster: the arbitrary time axis}
    There is no equivalent of the vertical ticks for the horizontal axis, so the
    raster branch does not attempt to calibrate time at this point. Instead the
    column indices are normalised to the unit interval,

    \[
      X_i = \frac{x_i - x_{\min}}{x_{\max} - x_{\min}} \in [0,1],
    \]
    and the axis is labelled \emph{arbitrary time scale}. The same $x_{\min}$ and
    $x_{\max}$ from the primary series are applied to the overlay and grow series,
    so all three share one time base and remain directly comparable.

    \paragraph{Orientation correction.} If the mean of the primary series is
    negative the sign of all three series is flipped. This handles displays that
    plot flow below the baseline, and applying it to all three series together
    preserves their alignment.

  \subsection{Raster: time-axis calibration from heart rate}
    Real time is recovered later, from the printed heart rate. The reasoning is
    that the mean interval between beat feet on the arbitrary axis must correspond
    to the true cardiac period $60/\mathrm{HR}$, giving a scale factor

    \[
      s = \frac{60/\mathrm{HR}}{T_{\text{arb}}},
      \qquad
      T_{\text{arb}} = \frac{x[f_{n-1}] - x[f_0]}{n_f - 1},
    \]
    where $f$ are the detected foot indices. \fn{digitized\_hr\_scale\_factor\_for\_raster}
    computes this from the ray series and the OCR dataframe, and
    \fn{scale\_raster\_digitized\_x} applies it to each series in turn. When the
    heart rate is missing or non-positive, or when foot spacing cannot be
    determined, the factor is $1.0$ and the axis remains arbitrary.\\
    \\
    This is the single most important dependency between stages: OCR failure on
    the heart rate line silently leaves the time axis uncalibrated. It also means
    the raster time axis inherits any error in the printed heart rate, which is
    why \fn{refine\_hr\_from\_local\_roi} is worth its extra OCR pass.

  \subsection{DICOM: mapping pixels using physical deltas}
    The DICOM branch needs no fitting because the display geometry is recorded in
    the header. \fn{\_dicom\_xy\_from\_curve\_coords} maps each condensed point
    directly:

    \begin{align*}
      x_{\text{ref}} &= \text{\fn{RegionLocationMinX0}} + \text{\fn{ReferencePixelX0}}, \\
      y_{\text{ref}} &= \text{\fn{RegionLocationMinY0}} + \text{\fn{ReferencePixelY0}}, \\
      X_i &= \text{\fn{ReferencePixelPhysicalValueX}}
             + (c_i - x_{\text{ref}})\cdot \delta_x, \\
      Y_i &= \text{\fn{ReferencePixelPhysicalValueY}}
             + (r_i - y_{\text{ref}})\cdot \delta_y .
    \end{align*}
    Here $\delta_x$ is \fn{PhysicalDeltaX} and
    $\delta_y = -\lvert \text{\fn{PhysicalDeltaY}} \rvert$. The reference pixel offsets
    default to zero, the reference physical values to zero and the deltas to one
    when absent. Crucially, the sign of $\delta_y$ is forced negative regardless
    of what the header states, because image rows
    increase downwards while velocity increases upwards. Consequently the DICOM
    branch has no need for the negative-mean inversion used on raster images, and
    does not apply it.\\
    \\
    \fn{Ynought} is set to \fn{[ReferencePixelPhysicalValueY]}, defaulting to
    $0.0$, whereas the raster branch always returns \fn{[0.0]}.\\
    \\
    The essential consequence is that DICOM $X$ is in physical units immediately,
    with no dependence on OCR, on beat detection or on the printed heart rate.
    Table~\ref{tab:calibcompare} summarises the difference.

    \begin{table}[h]
    \centering
    \small
    \begin{tabular}{@{}p{0.20\textwidth}p{0.36\textwidth}p{0.36\textwidth}@{}}
    \toprule
    & Raster & DICOM \\
    \midrule
    Velocity axis & Two-point fit through OCR-read ticks &
      Reference pixel plus \fn{PhysicalDeltaY} \\
    Time axis & Normalised to $[0,1]$, later scaled by heart rate &
      Reference pixel plus \fn{PhysicalDeltaX}, physical at once \\
    Requires OCR & Yes, for both ticks and heart rate & No \\
    \fn{Ynought} & \fn{[0.0]} & \fn{ReferencePixelPhysicalValueY} \\
    Sign handling & Inverted if mean is negative & \fn{PhysicalDeltaY} forced negative \\
    Failure mode & No ticks, so no calibration at all & Missing tags, so arithmetic error \\
    \bottomrule
    \end{tabular}
    \caption{Calibration compared across the two modalities.}
    \label{tab:calibcompare}
    \end{table}

  \subsection{Shape-preserving smoothing}
    Digitised traces retain single-pixel jitter, which would otherwise be
    amplified by the derivative-based foot detection of
    Section~\ref{sec:beats}. \fn{\_smooth\_1d\_digitized\_shape\_preserving}
    addresses this in two steps: a Hampel filter with a half-window of $2$ and a
    threshold of $3\sigma$ removes isolated spikes, then a Savitzky--Golay filter
    smooths the remainder. The window is $10\%$ of the series length clamped to
    $[7, 21]$ and forced odd, with polynomial order $3$ reduced if the window is
    too short. A rolling median is used as a fallback when the window would be
    smaller than five samples or the filter raises.\\
    \\
    A Savitzky--Golay filter is chosen over a moving average because it fits a
    local polynomial and therefore preserves peak height and curvature, which is
    essential when the peak values themselves are the measurements of interest.
    Smoothing is applied only to the ray series and to the grow series, since the
    morphological trace is already smooth by construction.

  \subsection{Primary and overlay series: an asymmetry between entry points}
  \label{sec:asymmetry}
    Both digitisation functions accept a primary curve and an overlay, with a flag
    \fn{overlay\_is\_ray} declaring which is which. The two entry points populate
    these arguments in opposite orders:

    \begin{itemize}
      \item \fn{single\_image\_processing.py} passes the \textbf{morphological}
        trace as primary with \fn{overlay\_is\_ray=True}.
      \item \fn{segment\_files.py} passes the \textbf{ray} trace as primary with
        \fn{overlay\_is\_ray=False}.
    \end{itemize}
    Plot colours and legends are driven by the flag and so remain correct in both
    cases. Metrics are also unaffected on the raster path, because
    \fn{plot\_correction} receives each series explicitly. On the DICOM path,
    however, \fn{waveform\_metrics\_from\_digitized} treats its primary
    \fn{(Xplot, Yplot)} argument as the ray series when labelling the output
    columns. The consequence is that for a DICOM file the column headed
    \emph{Digitized Value (ray)} contains morphological metrics when the
    single-image entry point is used, and ray metrics when the batch entry point
    is used. The numbers are correct in both cases; only the attribution differs.
    This is worth resolving before the column labels are relied upon
    quantitatively.
    %TODO: This probably isn't relevant at the moment since we only use this for debugging.

\newpage
\section{Stage 6: Beat Detection and Haemodynamic Indices}
\label{sec:beats}
  With a calibrated waveform in hand, the pipeline identifies individual cardiac
  cycles and reduces them to the standard haemodynamic indices. The beat detection
  logic is shared between modalities and lives in \fn{\_beat\_detection\_pass\_feet};
  only the wrapper differs.
  \fn{\_waveform\_peaks\_troughs\_values\_from\_physical\_x} is used when $x$ is
  already physical, as on the DICOM path, and
  \fn{\_digitized\_peaks\_metrics\_timescaled} when $x$ is arbitrary and must be
  scaled by heart rate, as on the raster path.

  \subsection{Why feet rather than troughs delimit a beat}
    A beat could in principle be delimited by successive minima, but the diastolic
    minimum of a Doppler envelope is a broad, flat feature whose location is
    sensitive to noise. The \emph{foot} — the onset of the systolic upstroke — is
    by contrast a sharp change in curvature and can be located far more
    repeatably. The pipeline therefore detects peaks first, uses them to anchor a
    search for feet, and only then finds the extrema within each foot-to-foot
    interval.

  \subsection{Peak proposal}
    Peaks are proposed with \fn{scipy.signal.find\_peaks} using two constraints
    derived from the signal itself rather than fixed thresholds:
    \begin{itemize}
      \item \textbf{Minimum separation.} A physiological upper bound of $200$\,bpm
        implies a minimum spacing of $0.3$\,s, converted to samples as
        $\max(1, \lfloor 0.3 / \Delta x \rfloor)$ where $\Delta x$ is the median
        sample spacing. Where $x$ is still arbitrary, the cruder
        $\max(1, n/15)$ is used, which assumes no more than fifteen beats on
        screen.
      \item \textbf{Prominence.} Set to a quarter of the peak-to-peak amplitude
        within \fn{\_beat\_detection\_pass\_feet}. Scaling with amplitude makes the
        criterion insensitive to the absolute velocity scale.
    \end{itemize}
    On the raster path this is done twice. The first pass runs on the arbitrary
    axis purely to obtain the mean foot spacing needed for the heart-rate scale
    factor; the second runs on the scaled axis with the physiological
    $0.3$\,s separation, and replaces the first result if it yields at least two
    feet and at least one peak and trough. The second pass cannot change the
    indices, since linear scaling of $x$ does not move samples, but it does change
    the separation constraint expressed in samples and therefore which peaks
    survive.

  \subsection{Foot detection}
    \fn{\_detect\_feet\_indices\_from\_peaks} locates one foot per anchoring peak
    using the derivatives of the envelope. For each peak:

    \begin{enumerate}
      \item Define a backward search window covering the last
        \fn{FOOT\_SEARCH\_FRACTION} $=0.50$ of the preceding peak-to-peak
        interval, tightened by $0.10$ for all beats after the first, and ending
        \fn{FOOT\_MIN\_SAMPLES\_BEFORE\_PEAK} $=3$ samples short of the peak. For
        the first peak the interval is estimated from the mean of the
        \emph{other} peak intervals, which is more stable than the first observed
        gap; the beat is skipped if the window would extend before the start of
        the signal.
      \item Smooth the window with a Savitzky--Golay filter of window at most
        \fn{FOOT\_DERIV\_SMOOTH\_WINDOW\_MAX} $=11$ and order
        \fn{FOOT\_DERIV\_SMOOTH\_POLYORDER} $=2$.
      \item Resample onto a uniform grid in $x$ before differentiating, so that
        irregular sample spacing cannot distort the gradients.
      \item Take the maximum of the second derivative as the upstroke anchor, the
        point of greatest positive curvature.
      \item Walk backwards from that anchor along the first derivative and accept
        the first low-slope point as the foot.
    \end{enumerate}

    A height guard rejects candidates that sit too high on the upstroke:
    \[
      y_{\text{foot}} \leq y_{\text{trough}}
        + \fn{FOOT\_MAX\_REL\_HEIGHT}\,(y_{\text{peak}} - y_{\text{trough}}),
      \qquad \fn{FOOT\_MAX\_REL\_HEIGHT} = 0.55 .
    \]
    Feet are returned sorted and deduplicated. Beat detection is abandoned, with
    an informational log entry, if fewer than two feet are found, since a single
    foot delimits no beat.

  \subsection{Per-beat extrema}
    \fn{\_peaks\_troughs\_from\_feet} treats beat $i$ as the half-open interval
    $[f_i, f_{i+1})$ and extracts one peak and one trough from each. A rolling
    median of window $5$ stabilises the selection. Where one of the originally
    proposed peaks falls inside the beat, the tallest such anchor is used in
    preference to the segment maximum, which keeps the reported peak consistent
    with the peak that defined the beat. The trough is then constrained to lie
    \emph{after} the peak, reflecting the physiological ordering of systole and
    diastole, and only falls back to the beat-wide minimum if the post-peak
    segment is empty. Intervals of three samples or fewer are skipped.

  \subsection{Signal quality filtering}
    \fn{\_sqi\_template\_correlation\_from\_feet} implements a template-matching
    quality index. Each beat is resampled to $200$ points, a median template is
    formed across all valid beats, and a beat is accepted if its Pearson
    correlation with the template is at least $0.85$. Beats that are degenerate —
    shorter than three samples, or with standard deviation below $10^{-10}$ — are
    marked invalid. Two fallbacks prevent the filter from discarding everything: if
    no beat is valid, or if the template itself is flat, or if every beat is
    rejected, all beats are accepted.

    \textbf{This filter is disabled by default.} \fn{USE\_SQI\_FILTER} is
    \fn{False} at \fn{general\_functions.py:276}, so every detected beat currently
    contributes to the reported metrics. The mechanism is complete and gated by a
    single flag, so enabling it is a one-line change, but the effect on reported
    values has not been characterised and it should be treated as experimental.

  \subsection{Index formulae}
    The indices themselves live in \fn{hemodynamic\_indices.py} as pure functions,
    deliberately separated from the image processing so that they can be tested
    and audited independently. Peak systolic and end-diastolic velocities are the
    means over all accepted beats,
    \[
      PS = \operatorname{mean}\big(y[\text{peaks}]\big), \qquad
      ED = \operatorname{mean}\big(y[\text{troughs}]\big),
    \]
    with $ED$ replaced by machine epsilon if it is exactly zero, guarding the
    divisions that follow. Let $\bar v = \operatorname{mean}(y)$ over
    \emph{all} envelope samples. Then:

    \begin{align*}
      S/D  &= PS / ED \\
      RI   &= \frac{PS - ED}{PS}
        && \text{\fn{resistive\_index\_from\_ps\_ed}} \\
      TA_{\max} &= \bar v
        && \text{\fn{tamax\_from\_envelope\_temporal\_mean}} \\
      PI   &= \frac{PS - ED}{\bar v}
        && \text{\fn{pulsatility\_index\_from\_ps\_ed\_and\_mean\_velocity}}
    \end{align*}

    All are rounded to two decimal places, and each guards against a zero or
    non-finite denominator by returning $0.0$.

    Two points deserve emphasis. First, $TA_{\max}$ is the temporal mean of the
    whole digitised envelope, not of complete beats only, so a partial beat at
    either end of the trace biases both $TA_{\max}$ and $PI$; the batch plots shade
    incomplete beats for exactly this reason. Second,
    \fn{tamax\_from\_ps\_ed\_approximation}, which computes
    $(PS + 2\,ED)/3$, is a legacy surrogate used only by the OCR consistency
    checks of Section~\ref{sec:ocr}, where no envelope is available. It is not used
    on any digitised path, and the two must not be conflated when comparing
    digitised and OCR values.

    Heart rate is not computed from the digitised trace. On the raster path it is
    an input, read by OCR and used to calibrate time; on the DICOM path it is
    neither needed nor reported.

\newpage
\section{Stage 7: Model Comparison and Selection}
\label{sec:selection}
  Stage 4 produces three candidate traces and the pipeline must return one. The
  criterion is agreement with the machine's own printed measurements: whichever
  trace reproduces the OCR values most closely is assumed to have segmented the
  envelope most faithfully. This makes the OCR stage not merely a reporting
  convenience but the arbiter of segmentation quality, and it is why an
  \fn{Fail} flag from Section~\ref{sec:ocr} degrades the result even though the
  segmentation itself succeeded.

  \emph{This stage applies to raster images only.} The DICOM path has no
  independent ground truth and is discussed at the end of the section.

  \subsection{Mapping digitised values into the metric table}
    \fn{plot\_correction} inserts three empty columns into the OCR dataframe at
    positions 3 to 5, named \emph{Digitized Value (ray)},
    \emph{Digitized Value (morph)} and \emph{Digitized Value (grow)}, so that they
    sit beside the extracted \texttt{Value} column for visual comparison.

    Metrics are then computed for each series — the ray series through
    \fn{plot\_correction} itself and the other two through
    \fn{\_digitized\_peaks\_metrics\_timescaled}, which reuses the primary
    series' arbitrary period so that all three share one time base. Each value is
    written into the row matched by \fn{\_df\_word\_metric\_mask}, which performs a
    substring match on the \texttt{Word} column with a dedicated regular
    expression for ophthalmic labels. The metric words are
    \texttt{PS, ED, S/D, RI, TA, PI} in general, and
    \texttt{PS, ED, PS/ED, RI, PI} for ophthalmic scans, which report no
    time-averaged velocity.

    The OCR \texttt{Value} column is never overwritten. It remains the reference
    against which the three digitised columns are scored.

  \subsection{Agreement scoring}
    \fn{digitized\_comparison.py} classifies each metric as a velocity or a ratio
    via \fn{metric\_kind\_for\_digitized\_row}, because an absolute error of
    $0.05$ means something very different for a velocity in centimetres per second
    than for a dimensionless index. Ratio metrics are those whose \texttt{Word}
    contains \texttt{PS/ED}, \texttt{ED/PS}, \texttt{S/D}, \texttt{PI} or
    \texttt{RI}; velocity metrics are those containing \texttt{TA}, \texttt{PS} or
    \texttt{ED}. The ratio tests are ordered before the bare \texttt{PS} and
    \texttt{ED} tests so that a compound label is not misclassified.

    Each cell is then assigned a colour. Let $t$ be the OCR value and $p$ the
    digitised value. For velocities, comparison is on magnitudes, since either may
    carry a spurious sign:
    \[
      AE = \big\lvert\, \lvert p \rvert - \lvert t \rvert \,\big\rvert,
      \qquad
      PE = 100\,\frac{AE}{\lvert t \rvert}\ \ (\lvert t \rvert > 0).
    \]

    \begin{table}[h]
    \centering
    \small
    \begin{tabular}{@{}lll@{}}
    \toprule
    Colour & Velocity rule & Ratio rule \\
    \midrule
    Green & $AE \leq 5$ \textbf{and} $PE \leq 5\%$        & $AE \leq 0.05$ \\
    Amber & otherwise                                      & $0.05 < AE \leq 0.10$ \\
    Red   & $AE > 10$ \textbf{or} $PE > 10\%$              & $AE > 0.10$ \\
    \bottomrule
    \end{tabular}
    \caption{Traffic-light thresholds as implemented in
      \fn{velocity\_digitized\_style} and \fn{ratio\_digitized\_style}.}
    \label{tab:trafficlight}
    \end{table}

    When $\lvert t \rvert$ is within $10^{-12}$ of zero the percentage error is
    undefined, and the cell is green only if $AE < 10^{-9}$ and red otherwise.

    \begin{warningbox}
    The docstrings in \fn{digitized\_comparison.py} do not match the code. The
    module docstring and \fn{velocity\_digitized\_style} state green at
    $AE \leq 5$ with $PE \leq 10\%$ and red at $AE > 15$ or $PE > 20\%$, whereas
    the implementation uses $PE \leq 5\%$ for green and $AE > 10$ or $PE > 10\%$
    for red. The implemented rules are therefore stricter than documented. The
    values in Table~\ref{tab:trafficlight} reflect the code. This should be
    reconciled, since the thresholds determine model selection and not merely cell
    colour.
    \end{warningbox}

  \subsection{Selecting the returned model}
    \fn{traffic\_light\_counts\_by\_metric\_kind} tallies each digitised column
    into two triples of green, amber and red counts, one for velocities and one
    for ratios, skipping any row whose OCR value or digitised value is missing,
    unparseable or non-finite. \fn{select\_best\_digitized\_model\_for\_image}
    then ranks the columns by the key

    \begin{lstlisting}[language=Python]
(velocity_greens, ratio_greens,
 -velocity_ambers, -ratio_ambers,
 -velocity_reds,   -ratio_reds)
    \end{lstlisting}

    sorted in descending order. Velocity greens dominate because peak systolic and
    end-diastolic velocities are the primary clinical measurements and the ratios
    are derived from them; a trace that gets the velocities right will generally
    get the ratios right, whereas the converse does not hold. Ambers are penalised
    before reds so that a column with several near misses is preferred to one that
    is confidently wrong. Two fallbacks apply: if no column scored any cell, the
    first column holding any numeric value is taken; and if no candidate column
    exists at all, \texttt{"ray"} is returned.

    \fn{finalize\_image\_digitized\_model\_choice} applies the decision. It
    consults the selector, then requires the chosen series to have at least two
    points with matching $x$ and $y$ lengths; if not, it falls back through ray,
    morph and grow in that order until a usable series is found. It records the
    outcome in a \texttt{Returned model} column as
    \fn{returned\_model\_cell\_value}, for example \texttt{"morph algorithm"}, and
    returns the selected $(x, y)$ pair. When no dataframe is supplied, or the
    dataframe lacks the ray column, it defaults to ray without scoring.

    The full raster ordering is therefore: digitise with
    \fn{plot\_digitized\_data\_single\_axis}; populate the digitised columns with
    \fn{plot\_correction}; compute and apply the heart-rate scale factor to all
    three series; and finally select with
    \fn{finalize\_image\_digitized\_model\_choice}.

  \subsection{DICOM: no independent ground truth}
    The DICOM path calls \fn{waveform\_metrics\_from\_digitized} instead of
    \fn{plot\_correction}, and neither
    \fn{select\_best\_digitized\_model\_for\_image} nor
    \fn{finalize\_image\_digitized\_model\_choice} is invoked. There is no
    \texttt{Returned model} column, and the trace returned is simply the primary
    series supplied by the caller — which, per Section~\ref{sec:asymmetry},
    differs between the two entry points.

    A structural consequence deserves attention. Because DICOM files carry no
    printed metric panel in the extracted region, the \texttt{Value} column of the
    resulting dataframe is populated with the computed primary metrics rather than
    with an independent measurement. The HTML report nonetheless applies the same
    traffic-light styling, comparing each digitised column against
    \texttt{Value}. The primary column therefore always scores perfect green, since
    it is being compared with itself, and the other two columns are scored against
    the primary rather than against any ground truth. The colours remain useful as
    an indication of agreement \emph{between} the three segmentation methods, but
    they must not be read as accuracy.

\newpage
\section{Batch Pipeline and Outputs}
\label{sec:batch}
  The batch entry point is \fn{usseg.main.main}, which composes three stages, each
  wrapped in the profiling helper \fn{prof} so that a
  \texttt{cProfile} report is written per stage:

  \begin{lstlisting}[language=Python]
usseg.setup_tesseract()
filenames = prof(usseg.get_likely_us, root_dir)
prof(usseg.segment, filenames)
prof(usseg.generate_html_from_pkl)
  \end{lstlisting}

  The stages communicate through pickle files rather than in memory, which means
  each can be re-run independently — a useful property when only the report needs
  regenerating after a long segmentation run.

  \subsection{File discovery: \fn{get\_likely\_us}}
    \fn{organise\_files.get\_likely\_us} walks \fn{root\_dir} with
    \fn{os.walk}, collecting files with extensions \fn{.jpg}, \fn{.png},
    \fn{.dcm} and \fn{.dicom}, and tests each with \fn{check\_file\_for\_us}.
    The two modalities are triaged differently:

    \begin{itemize}
      \item \textbf{DICOM} files are accepted unconditionally. The container
        format is itself sufficient evidence that the file is an ultrasound.
      \item \textbf{JPG and PNG} files are passed to \fn{scan\_type\_test}, which
        performs a cheap OCR probe and returns \fn{Fail} $=0$ for a probable
        Doppler scan. This is necessary because a directory of exported images
        typically also contains B-mode frames, report pages and screenshots.
    \end{itemize}

    Accepted files are grouped into a dictionary keyed by a patient identifier
    scraped from the file path with a regular expression, and that dictionary is
    pickled to \fn{likely\_us\_images}. Two details are worth recording. First,
    the identifier pattern differs by modality: \fn{\textbackslash d\{4\}} for
    DICOM and \fn{\textbackslash d\{3\}} for raster, so the same directory tree can
    yield inconsistent grouping. Second, the function returns a \emph{flattened}
    list of paths, discarding the grouping, which is preserved only in the pickle.
    Downstream code therefore processes scans independently, with no per-patient
    aggregation.

    An optional \fn{use\_parallel} flag routes discovery through a
    \fn{ThreadPoolExecutor}. It defaults to \fn{False} and \fn{main} does not
    enable it. Since the probe is dominated by Tesseract calls, which release the
    interpreter lock, threading would be expected to help; the conservative default
    presumably reflects Tesseract's thread-safety being uncertain in this
    configuration.

  \subsection{Per-scan processing: \fn{segment}}
    \fn{segment\_files.segment} accepts a list of paths, a single path, a
    dictionary, or a path to a pickle, normalising all of them to a list. Each
    scan is dispatched by extension to the DICOM or raster branch already described
    in Sections~\ref{sec:coarse} to~\ref{sec:selection}.

    \paragraph{Output naming.} Figures are written with the prefix
    \fn{output\_dir + f"\{idx\}\_\{base\_name\}"}. The loop index is included
    deliberately: DICOM exports frequently reuse basenames such as
    \fn{001.dcm} across patient folders, and without the index later scans would
    silently overwrite earlier ones.

    \paragraph{The parallel-list invariant.} The function accumulates five lists —
    \fn{scan\_display\_paths}, \fn{Digitized\_scans}, \fn{Annotated\_scans},
    \fn{Text\_data} and \fn{Mean\_wave\_scans} — which must remain index-aligned,
    because the report pairs them positionally with \fn{zip}. Every failure branch
    therefore appends \fn{None} rather than skipping, and this is the reason the
    exception handlers are structured as they are. A branch that returned early
    without appending would silently shift every subsequent row of the report
    against the wrong scan.

    \paragraph{Failure counting.} \fn{Fail} is incremented when digitisation
    raises. It is not aggregated into a run-level summary, so the count of failed
    scans must be recovered from the log or by counting \fn{None} entries in the
    pickle. A per-run summary would be a worthwhile addition.

    \paragraph{Memory hygiene.} After each scan the code calls
    \fn{plt.close("all")}, which is essential: the pipeline draws into figures 1
    and 2 by number, and without closing them the axes of successive scans would
    accumulate. A subsequent loop attempts to delete named locals through
    \fn{exec("del \%s" \% i)}. This has no effect in Python 3, because
    \fn{exec} cannot bind or unbind local variables in a function scope; it is dead
    code and can be removed.

  \subsection{Rendered outputs}
    For each scan the batch path may write four PNG files into
    \fn{output\_dir}:

    \begin{table}[h]
    \centering
    \small
    \begin{tabular}{@{}lp{0.55\textwidth}@{}}
    \toprule
    Suffix & Contents \\
    \midrule
    \fn{\_Source.png}     & DICOM only: the extracted Doppler frame, saved because a
                            browser cannot display raw \fn{.dcm} bytes \\
    \fn{\_Annotated.png}  & Figure 1: the original scan overlaid with the ROIs,
                            detected ticks and the three traces \\
    \fn{\_Digitized.png}  & Figure 2: the calibrated waveforms with beat markers,
                            saved at 900\,dpi \\
    \fn{\_MeanWaves.png}  & Mean beat per method, via
                            \fn{save\_mean\_waves\_ray\_morph\_grow\_figure} \\
    \bottomrule
    \end{tabular}
    \caption{Per-scan figures written by the batch pipeline.}
    \end{table}

    Figure 2 is saved by explicit number, \fn{plt.figure(2).savefig(...)}, because
    annotation may have left figure 1 current. The mean-wave figure is wrapped in
    its own handler and contributes \fn{None} on failure, so a mean-beat problem
    cannot cost the scan its other outputs.

  \subsection{Serialised artifacts}
    \fn{segment} pickles the five lists to \fn{segmented\_data} as a single
    list, in the order

    \begin{lstlisting}[language=Python]
[scan_display_paths, Digitized_scans, Annotated_scans,
 Text_data, Mean_wave_scans]
    \end{lstlisting}

    and returns the same data to its caller. Passing \fn{pickle\_path=False}
    suppresses the write, which is what the single-image and test paths do. Note
    that \fn{Text\_data} holds live \fn{pandas} dataframes, so the pickle is
    coupled to the pandas version used to write it.

  \subsection{HTML report}
    \fn{visualisation\_html.generate\_html\_from\_pkl} reads the pickle and
    calls \fn{generate\_html}, writing \fn{generated\_segmented\_data.html}
    into the working directory. It tolerates a four-element pickle from an earlier
    version by synthesising a list of \fn{None} for the mean-wave column.

    Each scan becomes one horizontally scrolling row of four images followed by its
    metric table, and every image is inlined as a base64 data URI so that the
    report is a single self-contained file. The trade-off is size: base64 inflates
    each image by roughly a third and nothing is shared between rows, so a report
    over a large batch becomes slow to open. Note also that the MIME type is
    hard-coded to \fn{image/png} even for JPEG inputs; browsers sniff the content
    and render it regardless, but the declaration is incorrect.

    Table cells in the three digitised columns are coloured by
    \fn{digitized\_cell\_background\_style} using the rules of
    Table~\ref{tab:trafficlight}, with empty cells left white. The styling call is
    wrapped in a bare \fn{except} that falls back to an unstyled cell, so a
    non-numeric value degrades the colour rather than the report. Column widths are
    estimated from the longest string in each column, which keeps the wide metric
    tables legible without a stylesheet.

\newpage
\section{Error Handling and Known Failure Modes}
\label{sec:failures}
  \subsection{Degradation strategy}
    The pipeline is designed to degrade rather than to abort. Each stage is
    wrapped in its own handler, and a failure removes only the information that
    stage would have contributed. The rationale is that a batch of several hundred
    scans will always contain some that fall outside the assumptions of the
    layout, and a single such scan must not cost the operator the whole run.

    The cascade of degradation is as follows:

    \begin{table}[h]
    \centering
    \small
    \begin{tabular}{@{}p{0.28\textwidth}p{0.30\textwidth}p{0.34\textwidth}@{}}
    \toprule
    Stage fails & Immediate effect & Downstream consequence \\
    \midrule
    One axis side & That side's ticks are \fn{None} &
      Calibration proceeds from the other side \\
    Both axis sides & No calibration possible &
      Empty digitised series; no metrics \\
    OCR heart rate & Scale factor is $1.0$ &
      Time axis stays arbitrary; velocities unaffected \\
    OCR overall (\fn{Fail}) & \fn{plot\_correction} skipped &
      No digitised columns, so model selection defaults to ray \\
    One tracing method & Its series is empty &
      Selection falls back through ray, morph, grow \\
    Digitisation & \fn{None} appended to all output lists &
      Report row shows blanks but stays aligned \\
    Mean-wave figure & \fn{None} in that column only &
      Other three figures still written \\
    \bottomrule
    \end{tabular}
    \caption{How failure at each stage propagates.}
    \end{table}

  \subsection{Commonly observed failures}
    Two exceptions dominate the console output of a healthy run and are worth
    recognising on sight, because both are expected rather than pathological.

    \paragraph{\texttt{ValueError: attempt to get argmax of an empty sequence}.}
    Raised inside \fn{search\_for\_ticks} when the column profile for one side has
    no peaks, that is when the scan prints no ticks on that side. Newer ophthalmic
    scans label the right axis only, so this fires once per image for the left
    side. It is caught, and the pipeline calibrates from the right axis alone.

    \paragraph{\texttt{ValueError: XA must be a 2-dimensional array}.}
    Raised by \fn{scipy.spatial.distance.cdist} inside
    \fn{search\_for\_labels} when there are no label centres to pair, which
    follows from the same cause. Also caught.

    Because \fn{segment\_files.py} calls \fn{traceback.print\_exc()} in these
    handlers as well as logging, both produce full tracebacks on stderr for every
    scan processed. The handling is correct but the reporting is misleading: a
    successful run looks like a failing one. Downgrading these two expected cases
    to a single-line debug message, while keeping full tracebacks for genuinely
    unexpected exceptions, would make the log usable as a monitoring tool.

    \paragraph{\texttt{FutureWarning} from \fn{binary\_erosion}.}
    \fn{\_shrink\_seed\_for\_region\_grow} calls
    \fn{morphology.binary\_erosion}, deprecated in scikit-image 0.26 and due for
    removal in 0.28. Behaviour is unchanged for now; the call should become
    \fn{morphology.erosion} with the same footprint.

  \subsection{Limitations}
    \begin{itemize}
      \item \textbf{Layout dependence.} Panel position, the yellow annotation
        colour and the tick-labelling convention are assumed throughout. The
        spatial pre-mask in \fn{text\_from\_greyscale} and the ROI fractions in
        \fn{extract\_dicom\_label\_text} are tuned to the machines in use and would
        need revisiting for another vendor.
      \item \textbf{Time calibration depends on OCR.} On raster images real time
        exists only if the printed heart rate is read correctly, and the pipeline
        cannot detect a plausible-but-wrong heart rate.
      \item \textbf{Single spectral region.} \fn{extract\_dicom\_metadata} takes
        the first matching pulsed-wave region and
        \fn{extract\_doppler\_from\_dicom} handles only single-frame
        $(H,W)$ and $(H,W,3)$ arrays, raising on multiframe data.
      \item \textbf{No independent validation on DICOM.} With no printed metric
        panel, the traffic-light colours measure inter-method agreement rather than
        accuracy, and no model selection is performed.
      \item \textbf{Endpoint-only axis fit.} A single surviving bad endpoint tick
        rescales every velocity on the scan.
      \item \textbf{Partial beats bias means.} $TA_{\max}$ and $PI$ average the
        whole envelope, including incomplete beats at either end.
      \item \textbf{Quality filtering inactive.} \fn{USE\_SQI\_FILTER} is
        \fn{False}, so no beat is excluded on quality grounds.
    \end{itemize}

\newpage
\section{Implementation Notes}
  \subsection{Performance considerations}
    Runtime is dominated by Tesseract. The threshold sweep in
    \fn{search\_for\_labels} may invoke OCR up to eighteen times per axis, and the
    early stop on divisibility by five exists primarily to cut that cost. Whole-scan
    OCR runs twice more in \fn{text\_from\_greyscale}, once through
    \fn{image\_to\_data} for boxes and once through \fn{image\_to\_string} for
    lines, plus a further constrained pass in \fn{refine\_hr\_from\_local\_roi}.

    \fn{colour\_extract\_vectorized} replaced a per-pixel implementation with
    NumPy array arithmetic and is no longer a significant cost;
    \fn{colour\_extract} remains in the module and is unused. Discovery can be
    threaded through \fn{use\_parallel}, though this is off by default.

    \fn{main} wraps each stage in \fn{prof}, writing a \texttt{cProfile} report
    named after the function, which is the intended starting point for any further
    optimisation.

  \subsection{Testing}
    Tests live in \fn{tests/} with fixture images in \fn{tests/resources}. Patient
    data is not committed, so \fn{tests/single\_image\_processing\_test.py}
    carries a placeholder path that must be pointed at a local scan; run directly,
    it also plots the extracted trace, which is the quickest way to sanity-check a
    change to the segmentation code.

  \subsection{Diagnostics and figure conventions}
    Figures are addressed by number: figure 1 is the annotated scan and figure 2
    the digitised waveforms. Traces are coloured consistently throughout —
    morphological in blue (\fn{\#1f77b4}), ray tracing in red, and region growing
    in violet (\fn{\#b026ff}) — and the same convention is used in this document.
    Several functions accept a \fn{debug\_plots} argument that renders intermediate
    masks; these are off by default because the batch path saves a fixed set of
    figures instead.

\newpage
\section{Recommended Follow-Up}
  The following items were identified while documenting the pipeline. They are
  listed in rough order of how much they affect interpretation of the results.

  \begin{enumerate}
    \item \textbf{Reconcile the traffic-light thresholds with their docstrings}
      (Section~\ref{sec:selection}). The implemented rules are stricter than
      documented and they drive model selection, so the discrepancy matters
      beyond presentation.
    \item \textbf{Make the primary/overlay order consistent between entry points}
      (Section~\ref{sec:asymmetry}), so that the \emph{Digitized Value (ray)}
      column means the same thing regardless of how a DICOM file was processed.
    \item \textbf{Quieten the two expected axis exceptions} so that a clean run
      produces a clean log, while preserving tracebacks for unexpected failures.
    \item \textbf{Replace \fn{morphology.binary\_erosion}} with
      \fn{morphology.erosion} before scikit-image 0.28 removes it.
    \item \textbf{Characterise and then enable the SQI filter}, or remove it, so
      that its status is unambiguous.
    \item \textbf{Restrict $TA_{\max}$ and $PI$ to complete beats}, since the
      incomplete-beat shading already identifies the samples to exclude.
    \item \textbf{Enforce the monotonicity result in \fn{validate\_axis\_ticks}},
      which is currently computed and discarded.
    \item \textbf{Unify the patient-identifier regular expressions} in
      \fn{check\_file\_for\_us}, presently four digits for DICOM and three for
      raster.
    \item \textbf{Report a per-run summary} of processed, degraded and failed
      scans, which the parallel-list structure already makes easy to compute.
    \item \textbf{Remove dead code}: the \fn{exec("del ...")} loop in
      \fn{segment}, the unused \fn{colour\_extract}, and the unread
      \fn{patient\_paths} configuration key.
  \end{enumerate}

\end{document}
